\section{Conclusion}
\label{sec:conclusion}

We presented OPENER, an open-world NER pipeline assembled entirely from ready-to-use components, a frozen detector, a pretrained Matryoshka embedder sharpened by a light contrastive fine-tuning with error-driven hard-negative mining, and a swappable typing head, with no encoder trained from scratch.
The modular design is part of the result rather than an implementation detail. Because the three stages are independently swappable, we could search the architecture one axis at a time (\autoref{tab:ablation}) and select every component on measured evidence instead of by assumption, and the same modularity keeps the embedder's Matryoshka structure usable as a compute/quality knob even after our contrastive fine-tuning (\autoref{fig:matryoshka}).

Both operating points of the typing head perform strongly, at opposite ends of the supervision axis.
Our zero-shot head OPENER-ZS, label-name prototypes refined transductively and fused with the detector's own predictions, needs no target labels at typing time and is still the best of the compared zero-shot systems end-to-end ($39.5$ AMI, all AMI scores being reported $\times 100$ as in the tables), ahead of GLiNER-L, GNER, Qwen2.5 and OWNER, at one to two orders of magnitude lower energy than the trained-encoder and generative baselines.
When a domain's entity classes are fixed and labelled, a balanced linear probe, OPENER-Sup, goes further. On gold mentions, the contrastive space is highly discriminative, and OPENER-Sup is the most accurate system in our benchmark, leading on all $13$ sets and surpassing a zero-shot transfer of OWNER ($62.3$ vs $43.0$ AMI). It also attains the best end-to-end AMI of the compared systems ($40.1$ AMI*100).

Both results rest on a benchmark that is itself a contribution: $13$ open-world NER datasets on which quality, per-sentence latency and energy are reported jointly, typing-on-gold is kept separate from end-to-end, and the supervision asymmetry between systems is stated rather than hidden. Every OPENER number is averaged over three full embedder-retraining seeds.

Two findings carry beyond the specific system.
First, an open-world typing head need not train its own encoder: a general-purpose embedding, once sharpened by a light contrastive step, is already discriminative enough to type entities, and the cost that does matter is fitting the cheap typing head, not the embedder.
Second, the three-axis evaluation locates the open-world bottleneck in \emph{detection} rather than typing: typing on gold mentions reaches $62.3$ AMI but end-to-end scores drop to $40.1$ as detector recall falls on specialised domains. So, improving the detector, not the typing head, is the clearest axis for future gains. The evaluation also exposes efficiency gaps that an accuracy-only comparison would hide.

The approach has clear limits.
End-to-end quality is capped by detector recall on specialised domains, the shared bottleneck, not by the typing head.
Supervised typing needs target-domain labels, and while OPENER-ZS removes them at typing time it depends on the quality of the label names and, in its transductive form, on access to the unlabelled test mentions. The inductive zero-shot head is correspondingly weaker.
Finally, our experiments run on deliberately modest hardware: a single consumer GPU with $6$\,GB of VRAM, subject to thermal throttling.
This is the very setting that makes OPENER deployable, but it also bounded our experimental budget: despite an already broad benchmark ($13$ datasets, three axes and extensive ablations), we cap each test split at $1000$ sentences.
The quality scores themselves are nonetheless well characterised: retraining the full embedder under three seeds leaves the benchmark means nearly unchanged (std of $0.1$ to $1.3$ AMI on the thirteen-set means, \autoref{sec:exp:protocol}), inference is deterministic at a fixed seed, and a test-set bootstrap confirms per-dataset stability (bootstrap std of at most $1.5$ AMI, \autoref{tab:bootstrap}). Only latency and energy fluctuate ($\pm 15\%$) and are therefore read as averages rather than precise points.

Future work includes improving mention detection on specialised domains to lift the end-to-end ceiling and extending the decoupled detect-then-type design beyond NER to entity linking and relation extraction.
